% ===================================================================== PREAMBULE COMMUN — Carnet d'entraînement Fiche 9
\documentclass[11pt,a4paper]{article}
\usepackage[utf8]{inputenc}
\usepackage[T1]{fontenc}
\usepackage{lmodern}
\usepackage[french]{babel}
\usepackage{amsmath,amssymb,mathtools}
\usepackage{icomma}
\usepackage{siunitx}
\sisetup{output-decimal-marker={,},per-mode=symbol}
\DeclareSIUnit{\molar}{M}
\DeclareSIUnit{\Molar}{mol\,L^{-1}}
\usepackage[version=4]{mhchem}
\usepackage{xcolor}
\usepackage{colortbl}
\usepackage{tikz}
\usepackage{pgfplots}
\usepackage[most]{tcolorbox}
\usepackage{enumitem}
\usepackage{array}
\usepackage{booktabs}
\usepackage{multicol}
\usepackage{fancyhdr}
\usepackage[hidelinks]{hyperref}
\usetikzlibrary{arrows.meta,positioning,calc,decorations.pathmorphing,
  decorations.markings,angles,quotes,patterns,backgrounds,
  shapes.geometric,shapes.misc,fadings,intersections,fit}
\pgfplotsset{compat=1.18}
\usepackage[a4paper,margin=2.2cm,headheight=26pt,headsep=10pt]{geometry}

% ---------- Palette ----------
\definecolor{primary}{RGB}{150,98,162}
\definecolor{primarydark}{RGB}{92,52,110}
\definecolor{defcol}{RGB}{0,114,178}
\definecolor{thmcol}{RGB}{0,140,120}
\definecolor{formulacol}{RGB}{198,93,9}
\definecolor{remarkcol}{RGB}{120,94,200}
\definecolor{attncol}{RGB}{200,40,55}
\definecolor{excol}{RGB}{0,135,90}
\definecolor{methcol}{RGB}{90,90,100}
\definecolor{cationcol}{RGB}{0,90,190}
\definecolor{anioncol}{RGB}{200,60,55}
\definecolor{cristalcol}{RGB}{110,105,120}
\definecolor{eaucol}{RGB}{120,180,220}
\definecolor{solcol}{RGB}{0,135,90}
\definecolor{kpscol}{RGB}{198,93,9}
\definecolor{qcol}{RGB}{120,94,200}
\definecolor{precipcol}{RGB}{110,70,40}
\definecolor{exercol}{RGB}{150,98,162}
\definecolor{corrcol}{RGB}{0,120,100}

% ---------- Macros ----------
\newcommand{\Kps}{K_{\mathrm{ps}}}
\newcommand{\Ce}[1]{\left[\,\ce{#1}\,\right]}

% ---------- Style des boîtes ----------
\tcbset{boxbase/.style={enhanced,breakable,boxrule=0.9pt,arc=2.5pt,
   left=3mm,right=3mm,top=2.3mm,bottom=2.3mm,
   fonttitle=\bfseries,coltitle=white,toptitle=0.6mm,bottomtitle=0.6mm}}

\newtcolorbox[auto counter]{definition}[1][]{%
  boxbase,colback=defcol!5,colframe=defcol,
  title={\textsc{Définition}~\thetcbcounter\ \ifx\relax#1\relax\relax\else\textmd{--- \itshape #1}\fi}}
\newtcolorbox[auto counter]{theoreme}[1][]{%
  boxbase,colback=thmcol!6,colframe=thmcol,
  title={\textsc{Loi / Propriété}~\thetcbcounter\ \ifx\relax#1\relax\relax\else\textmd{--- \itshape #1}\fi}}
\newtcolorbox[auto counter]{formule}[1][]{%
  boxbase,colback=formulacol!6,colframe=formulacol,
  title={\textsc{Formule}~\thetcbcounter\ \ifx\relax#1\relax\relax\else\textmd{--- \itshape #1}\fi}}
\newtcolorbox{remarque}[1][]{%
  boxbase,colback=remarkcol!6,colframe=remarkcol,
  title={\textsc{Remarque}\ifx\relax#1\relax\relax\else\textmd{ : \itshape #1}\fi}}
\newtcolorbox{attention}[1][]{%
  boxbase,colback=attncol!6,colframe=attncol,
  title={\textsc{Attention}\ifx\relax#1\relax\relax\else\textmd{ : \itshape #1}\fi}}
\newtcolorbox{exemple}[1][]{%
  boxbase,colback=excol!5,colframe=excol,
  title={\textsc{Exemple}\ifx\relax#1\relax\relax\else\textmd{ : \itshape #1}\fi}}
\newtcolorbox{methode}[1][]{%
  boxbase,colback=methcol!7,colframe=methcol,sharp corners,
  title={\textsc{Méthode}\ifx\relax#1\relax\relax\else\textmd{ : \itshape #1}\fi}}
\newtcolorbox{astuce}[1][]{%
  boxbase,colback=primary!7,colframe=primary,
  title={\textsc{Astuce}\ifx\relax#1\relax\relax\else\textmd{ : \itshape #1}\fi}}

\newtcolorbox[auto counter]{exercice}[1][]{%
  boxbase,colback=exercol!4,colframe=exercol,
  title={\textsc{Exercice}~\thetcbcounter\ \ifx\relax#1\relax\relax\else\textmd{--- \itshape #1}\fi}}
\newtcolorbox[auto counter]{correction}[1][]{%
  boxbase,colback=corrcol!4,colframe=corrcol,
  title={\textsc{Corrigé}~\thetcbcounter\ \ifx\relax#1\relax\relax\else\textmd{--- \itshape #1}\fi}}

% ---------- Environnement « où : » ----------
\newenvironment{ou}{%
  \par\smallskip\noindent\textbf{où :}\nopagebreak
  \begin{itemize}[leftmargin=1.8em,topsep=2pt,itemsep=1.5pt,parsep=0pt]}%
  {\end{itemize}}

% ---------- Styles TikZ ----------
\tikzset{
  cat/.style={circle,fill=cationcol,draw=cationcol!50!black,inner sep=0pt,
              minimum size=5mm,font=\bfseries\scriptsize,text=white},
  an/.style={circle,fill=anioncol,draw=anioncol!50!black,inner sep=0pt,
             minimum size=5mm,font=\bfseries\scriptsize,text=white},
  crx/.style={circle,fill=cristalcol,draw=black!50,inner sep=0pt,
              minimum size=5mm,font=\bfseries\scriptsize,text=white},
  ptn/.style={circle,fill=black,inner sep=1.1pt}}

% ---------- En-têtes ----------
\pagestyle{fancy}
\fancyhf{}
\fancyhead[L]{\small\textcolor{primarydark}{\textbf{Fiche 9} — Solubilité et produit de solubilité}}
\fancyhead[R]{\small\textcolor{primarydark}{\textsc{Carnet d'entraînement}}}
\fancyfoot[C]{\small\thepage}
\renewcommand{\headrulewidth}{0.8pt}
\renewcommand{\headrule}{\hbox to\headwidth{\color{primary}\leaders\hrule height \headrulewidth\hfill}}
\usepackage{titlesec}
\titleformat{\section}{\Large\bfseries\color{primarydark}}{\thesection}{0.6em}{}
\titleformat{\subsection}{\large\bfseries\color{primary}}{\thesubsection}{0.6em}{}

% ---------- Bandeau de titre ----------
% #1 = thème/étiquette   #2 = titre principal   #3 = sous-titre
\newcommand{\bandeau}[3]{%
\begin{center}
\begin{tikzpicture}
\node[rectangle,rounded corners=7pt,fill=primary,text=white,
      minimum width=15.6cm,minimum height=1.95cm,align=center,inner sep=8pt]
  {{\large\bfseries #1}\\[3pt]{\Large\bfseries #2}\\[3pt]{\small #3}};
\end{tikzpicture}
\end{center}
\vspace{2mm}}
\begin{document}
\bandeau{Thème 1 — Équation de dissolution \& $\Kps$}%
{Corrigé — Niveau Facile}%
{Réponses détaillées et justifications}

\begin{correction}[équations de dissolution simples]
On place le solide $(s)$ à gauche, les ions $(aq)$ à droite, on équilibre les charges.
\begin{enumerate}[label=\alph*),leftmargin=2em,topsep=2pt,itemsep=2pt]
  \item $\ce{AgCl(s) <=> Ag+(aq) + Cl-(aq)}$ \quad (type $1{:}1$).
  \item $\ce{CaCO3(s) <=> Ca^{2+}(aq) + CO3^{2-}(aq)}$ \quad (le groupement \ce{CO3^{2-}} reste entier ; type $1{:}1$).
  \item $\ce{KI(s) <=> K+(aq) + I-(aq)}$ \quad (type $1{:}1$).
\end{enumerate}
Dans les trois cas, une unité de sel donne \textbf{un} cation et \textbf{un} anion : aucun coefficient
supérieur à 1.
\end{correction}

\begin{correction}[type $p{:}q$]
\begin{enumerate}[label=\alph*),leftmargin=2em,topsep=2pt,itemsep=2pt]
  \item $\ce{PbBr2(s) <=> Pb^{2+}(aq) + 2 Br-(aq)}$ : 1 cation, 2 anions $\Rightarrow$ type $\boldsymbol{1{:}2}$.
  \item $\ce{CaF2(s) <=> Ca^{2+}(aq) + 2 F-(aq)}$ : 1 cation, 2 anions $\Rightarrow$ type $\boldsymbol{1{:}2}$.
  \item $\ce{Fe(OH)3(s) <=> Fe^{3+}(aq) + 3 OH-(aq)}$ : 1 cation, 3 anions $\Rightarrow$ type $\boldsymbol{1{:}3}$.
\end{enumerate}
\textbf{Contrôle des charges} pour (a) : à droite $(+2)+2\times(-1)=0$ ; à gauche le solide est neutre.
OK. Même vérification pour (b) et (c).
\end{correction}

\begin{correction}[règle d'or sur \ce{Ag3PO4}]
On lit directement les coefficients de l'équation : \textbf{3} devant \ce{Ag+}, \textbf{1} devant
\ce{PO4^{3-}}. Ces coefficients deviennent les exposants :
\[ \boxed{\ \Kps = \Ce{Ag+}^{3}\,\Ce{PO4^3-}\ } \]
\begin{itemize}[leftmargin=1.5em,topsep=1pt,itemsep=1pt]
  \item exposant \textbf{3} sur \ce{Ag+} : c'est le coefficient 3 de l'équation ;
  \item exposant \textbf{1} (implicite) sur \ce{PO4^{3-}} : coefficient 1.
\end{itemize}
Le solide \ce{Ag3PO4(s)} n'apparaît pas (activité $=1$). C'est la réponse du \textsc{qcm}\,n\textsuperscript{o}\,1
du livre : $\Kps=\Ce{Ag+}^{3}\Ce{PO4^3-}$.
\end{correction}

\begin{correction}[$\Kps$ à partir d'équations données]
\begin{enumerate}[label=\alph*),leftmargin=2em,topsep=2pt,itemsep=2pt]
  \item $\Kps=\Ce{Mn^2+}\,\Ce{OH-}^{2}$ \quad (exposant 2 sur \ce{OH-}, coefficient 2).
  \item $\Kps=\Ce{Ag+}^{2}\,\Ce{CrO4^2-}$ \quad (exposant 2 sur \ce{Ag+}).
  \item $\Kps=\Ce{Al^3+}\,\Ce{OH-}^{3}$ \quad (exposant 3 sur \ce{OH-}).
\end{enumerate}
Règle unique : \emph{exposant = coefficient stœchiométrique}. Le cation et l'anion apparaissent, le
solide jamais.
\end{correction}

\begin{correction}[nature du $\Kps$ : vrai/faux]
\begin{enumerate}[label=\alph*),leftmargin=2em,topsep=2pt,itemsep=2pt]
  \item \textbf{Faux.} Le solide pur a une activité de 1 : il est absent du $\Kps$. Seuls les ions
        \emph{dissous} y figurent.
  \item \textbf{Faux.} Le $\Kps$ est \textbf{sans unité} (constante d'équilibre). C'est la
        \emph{solubilité} $s$, elle, qui s'exprime en \si{\mole\per\litre}.
  \item \textbf{Vrai.} Par définition, $\Kps$ est la constante d'équilibre de la réaction de dissolution.
  \item \textbf{Vrai.} Comme toute constante d'équilibre, $\Kps$ dépend de $T$ (souvent il augmente
        avec $T$ car la dissolution est le plus souvent endothermique).
\end{enumerate}
\emph{(Ce sont les items C et D corrects du \textsc{qcm}\,n\textsuperscript{o}\,12 du livre.)}
\end{correction}

\end{document}
